\scanpage{169}
% archive foliation: 168
\uline{Application} Un revêt. étale \inserted{$T$ sur $S$} est dit à \uline{groupe structural} $\pi$, \inserted{\fbox{fibre $F$}}, \deleted{\unsure{si on a fixé un} sous-schéma de $\operatorname{Aut}_S(T)$, qui est} \unsure{si on a} donné une opération fidèle de $\pi$ sur $F$, et un sous-fibré \inserted{$\pi_S$} principal de $\underline{\operatorname{Isom}}_S(E_S,T)$. Si $S$ est connexe et muni d'un pt géométrique $\xi$, il revient au $=$ de se donner un $T$ avec une telle structure, ou de se donner des opérations \deleted{continues} de $\widetilde\pi_1(S,\xi)$ \deleted{\unsure{à travers}} sur $F$ :
\[
\text{hom}\qquad\widetilde\pi_1(S,\xi)\longrightarrow\deleted{\unsure{G}}\pi.
\]

\uline{Exemples} Schémas en groupes comm., constants tordus sur $S$, \deleted{à fibre \unsure{de}} \unsure{génération} finie sur $S$, correspondent \deleted{aux} $(\widetilde\pi_1(S,\xi))$-Modules admissibles \deleted{de \unsure{présent}} qui sont des $\mathbf Z$-modules de type fini.
